% =============================================================================
% 纸笔外语学习 - LaTeX 宏包定义
% Paper-and-Pencil Language Learning - Macro Definitions
% =============================================================================

% CEFR 等级标记
\providecommand{\cefr}[1]{\textsuperscript{\textsf{CEFR #1}}}
\newcommand{\cefrAone}{\textsuperscript{\textsf{CEFR A1}}}
\newcommand{\cefrAtwo}{\textsuperscript{\textsf{CEFR A2}}}
\newcommand{\cefrBone}{\textsuperscript{\textsf{CEFR B1}}}
\newcommand{\cefrBtwo}{\textsuperscript{\textsf{CEFR B2}}}
\newcommand{\cefrCone}{\textsuperscript{\textsf{CEFR C1}}}
\newcommand{\cefrCtwo}{\textsuperscript{\textsf{CEFR C2}}}

% 体裁标记
\providecommand{\genre}[1]{\textbf{[体裁：#1]}}
\newcommand{\genrenarrative}{\genre{记叙文}}
\newcommand{\genreexposition}{\genre{说明文}}
\newcommand{\genreargumentation}{\genre{议论文}}
\newcommand{\genrecorrespondence}{\genre{书信体}}

% 字数统计
\providecommand{\wordcount}[1]{\textbf{[字数要求：#1]}}
\newcommand{\wordcountrange}[2]{\textbf{[字数要求：#1-#2 词]}}

% 难度分级
\providecommand{\difficulty}[1]{\textcolor{blue!70!black}{\small\textbf{难度：}#1}}

% IPA 音标标注（保留，用于拼读训练）
\newcommand{\ipa}[1]{\textsf{/ #1 /}}
\newcommand{\ipaeth}{\ipa{θ}}
\newcommand{\ipath}{\ipa{ð}}
\newcommand{\iparl}{\ipa{l}}
\newcommand{\iparr}{\ipa{r}}

% 填空格式
\newcommand{\fillanswer}[2][2cm]{\underline{\hspace{#1}}\studentonly{\textsuperscript{\tiny[#2]}}}
\newcommand{\shortfill}[1][1cm]{\fillanswer[#1]}
\newcommand{\longfill}[1][4cm]{\fillanswer[#1]}

% 错误标注与纠正
\providecommand{\error}[2]{\uline{\sout{#1}}\teacheronly{\textcolor{red!70!black}{\textbf{→ #2}}}}
\newcommand{\correct}[1]{\textcolor{green!60!black}{\textbf{✓ #1}}}

% 使用短语框
\providecommand{\usefulphrase}[1]{%
  \vspace{0.3em}%
  \noindent\framebox[\linewidth][l]{%
    \parbox[t]{\dimexpr\linewidth-2\fboxsep}[t]{%
      \small\textbf{💡 Useful phrase:} #1%
    }%
  }%
  \vspace{0.3em}%
}

% 交际功能标记
\newcommand{\speechact}[1]{\textbf{[功能：#1]}}
\newcommand{\req}{\speechact{请求}}
\newcommand{\sug}{\speechact{建议}}
\newcommand{\apolo}{\speechact{道歉}}
\newcommand{\thank}{\speechact{感谢}}
\newcommand{\complain}{\speechact{投诉}}
\newcommand{\recommend}{\speechact{推荐}}

% 真实语料标记
\newcommand{\auth}[1]{\textcolor{blue!60!black}{\small\textit{[\textasciitilde 语料来源：#1]}}}
\newcommand{\authcoca}{\auth{COCA}}
\newcommand{\authbnc}{\auth{BNC}}

% 翻译对照
\newcommand{\chinesetag}[1]{\textbf{【中文】}#1}
\newcommand{\englishtag}[1]{\textbf{[英文]} #1}

% =============================================================================
% v2.0: 翻译训练相关宏（新增）
% =============================================================================

% 中文原文输入
\newcommand{\chineseinput}[1]{%
  \begin{chinesetext}
  #1
  \end{chinesetext}%
}

% 语法焦点提示
\newcommand{\transfocus}[1]{%
  \syntaxfocus{#1}%
}

% 字数要求简化版
\newcommand{\transwordlimit}[1]{%
  \wordlimit{#1}%
}

% 关键词汇对照条目
\newcommand{\transkeyword}[2]{%
  \item #1: #2
}

% 评分维度标签（v2.0 翻译标准）
\newcommand{\criterionfidelity}{\textbf{内容忠实度 (40\%)}}
\newcommand{\criteriongrammaraccuracy}{\textbf{语法准确性 (30\%)}}
\newcommand{\criterionlexicalappropriateness}{\textbf{词汇地道性 (20\%)}}
\newcommand{\criterioncohesion}{\textbf{衔接连贯性 (10\%)}}

% 常见翻译错误预警
\newcommand{\translationerrorwarning}[2]{%
  \teacheronly{%
    \vspace{0.3em}%
    \framebox[\linewidth][l]{%
      \parbox[t]{\dimexpr\linewidth-2\fboxsep}[t]{%
        \small\textbf{⚠ 常见中式英语：}"#1" \\- \\_\textbf{应译为：}"#2"
      }%
    }%
    \vspace{0.3em}%
  }%
}

% 跨文化翻译注释
\newcommand{\culturetranslatip}[2]{%
  \marginpar{%
    \framebox[\linewidth][l]{%
      \parbox[t]{\marginparwidth-2\fboxsep}{%
        \small\textbf{🌍 文化翻译提示}\\
        #1 \\- \\#2
      }%
    }%
  }%
}

% =============================================================================
% v2.0: 移除写作任务标识（已废弃）
% 说明：原 writingtasktype/writingaudience/writingtone 等命令不再使用
%       保留向后兼容但标记为 deprecated
% =============================================================================

% 旧版写作命令（deprecated after v2.0）
% \newcommand{\writingtasktype}[1]{\textbf{任务类型：}#1}  % v2.0 DEPRECATED
% \newcommand{\writingaudience}[1]{\textbf{读者对象：}#1}  % v2.0 DEPRECATED
% \newcommand{\writingtone}[1]{\textbf{语气风格：}#1}      % v2.0 DEPRECATED

% 评分维度标签
\newcommand{\criterioncontent}{\textbf{内容 (35\%)}}
\newcommand{\criterionlanguage}{\textbf{语言 (35\%)}}
\newcommand{\criterionorg}{\textbf{组织 (20\%)}}
\newcommand{\criterionregister}{\textbf{得体性 (10\%)}}

% 段落结构提示
\newcommand{\introductoryparagraph}{\textbf{引言段}}
\newcommand{\bodyparagraph}[1]{\textbf{主体段 #1}}
\newcommand{\concludingparagraph}{\textbf{结论段}}

% 衔接手段识别
\newcommand{\transition}[1]{\textcolor{purple!70!black}{\textbf{[连接词：#1]}}}
\newcommand{\reference}[1]{\textcolor{orange!70!black}{\textbf{[指代：#1]}}}
\newcommand{\ellipsismark}[1]{\textcolor{green!60!black}{\textbf{[省略：#1]}}}

% 词汇链追踪
\newcommand{\lexicalchain}[1]{\textcolor{teal!70!black}{\textbf{[词汇链：#1]}}}

% 常见错误预警
\newcommand{\commonerrorwarning}[2]{%
  \teacheronly{%
    \vspace{0.3em}%
    \framebox[\linewidth][l]{%
      \parbox[t]{\dimexpr\linewidth-2\fboxsep}[t]{%
        \small\textbf{⚠ 学生常犯错误：} "#1" \quad $\rightarrow$ \textbf{应为：} "#2"%
      }%
    }%
    \vspace{0.3em}%
  }%
}

% 跨文化注释
\newcommand{\culturetip}[1]{%
  \marginpar{%
    \framebox[\linewidth][l]{%
      \parbox[t]{\marginparwidth-2\fboxsep}{%
        \small\textbf{🌍 文化提示}\\#1%
      }%
    }%
  }%
}

% 学习策略提示
\newcommand{\learningtip}[1]{%
  \vspace{0.3em}%
  \noindent\framebox[\linewidth][l]{%
    \parbox[t]{\dimexpr\linewidth-2\fboxsep}[t]{%
      \small\textbf{💡 学习小贴士：} #1%
    }%
  }%
  \vspace{0.3em}%
}

% 自我评估工具
\newcommand{\selfcheckstart}{%
  \vspace{0.5em}%
  \textbf{自我评估：}%
  \begin{itemize}[leftmargin=*, label=$\square$]%
}
\newcommand{\selfcheckend}{\end{itemize}\vspace{0.3em}}

% 同伴互评提示
\newcommand{\peerreviewprompt}[1]{%
  \vspace{0.3em}%
  \textbf{同伴互评：}#1%
}

% 范文批注标记
\newcommand{\modelnote}[1]{%
  \teacheronly{%
    \textcolor{blue!70!black}{\textbf{【注释：}#1\textbf{]}}%
  }%
}

% 梯度练习标识（v2.0 翻译层级）
\newcommand{\controlledpractice}{\textbf{Level 1: 单句精准翻译}}
\newcommand{\guidedpractice}{\textbf{Level 2: 段落要点翻译（强支持）}}
\newcommand{\freeproduction}{\textbf{Level 3: 观点转换表达（中等支持）}}
% Note: v2.0 removed Level 4 independent writing

% 任务步骤编号
\newcommand{\langstep}[2]{\textbf{步骤 #1:} #2}
\providecommand{\step}{\langstep}

% 时间建议
\newcommand{\estimatedtime}[1]{\textbf{建议用时：}#1}

% 答案揭示控制
\newcommand{\revealanswer}[1]{%
  \studentonly{}%
  \teacheronly{\vspace{0.3em}\textbf{参考答案：}#1}%
}

% 多版本答案提示
\newcommand{\multipleanswersacceptable}{%
  \teacheronly{%
    \vspace{0.3em}%
    \textbf{【注】本题答案不唯一，合理即可}%
  }%
}

% 变体答案支持
\providecommand{\alternativeanswer}[1]{%
  \teacheronly{%
    \item \textit{替代答案：}#1%
  }%
}

% 部分得分标准
\newcommand{\partialcredit}[2]{%
  \teacheronly{%
    \vspace{0.2em}%
    \small\textcolor{blue!70!black}{#1 \quad ($\rightarrow$ #2分)}%
  }%
}

% 典型回答示例
\newcommand{\sampleanswer}[1]{%
  \teacheronly{%
    \vspace{0.3em}%
    \textit{【示例】}#1%
  }%
}

% 认知负荷提示
\newcommand{\cognitiveload}[1]{%
  \teacheronly{%
    \vspace{0.2em}%
    \small\textcolor{red!70!black}{\textbf{认知负荷：}#1}%
  }%
}

%  scaffolding level
\newcommand{\scaffoldinglevel}[1]{%
  \teacheronly{%
    \vspace{0.2em}%
    \small\textcolor{green!70!black}{\textbf{脚手架层级：}#1}%
  }%
}

% 文本复杂度分析
\newcommand{\textcomplexityanalysis}[3]{%
  \teacheronly{%
    \vspace{0.3em}%
    \small\textbf{文本复杂度分析：}%
    \begin{itemize}[leftmargin=*, nosep, label=$\bullet$]%
      \item 平均句长：#1 词
      \item 生词密度：#2%
      \item CEFR 对标：#3%
    \end{itemize}%
    \vspace{0.3em}%
  }%
}
